\section{Problem statement}
\subsection{Context}
\noindent
Catalysis contributes to \textbf{a third} of the global economy.
It dictates our lives through driving energy and food security, pharmaceutical synthesis, and manufacturing.
The theoretical exploration of new catalysts is required to further accelerate development in these key areas.
Theoretical exploration expedites catalyst discovery due to removing the need to experimentally synthesize and test candidates.

\subsection{Problem}
\noindent


\subsection{Significance}
\noindent


\subsection{Goals}
\noindent
This project will be deemed a success when:
\begin{itemize}
    \item The online KMC model reproduces results of other KMC models.
    \item The product and dominant reaction pathway is found for a process, with only reactants and surface structure given as an input.
    \item Kinetic modeling is performed for a system that is unfeasible with current methods.
\end{itemize}

\section{Theory}
\subsection{What is a bond?}
\noindent
Bonding can be separated into two categories: ionic and covalent.
Ionic bonding can be explained simply from an electron transfer from one species to another, forming charged atoms.
These charged species are attracted to each other based on electrostatic attraction.
For covalent bonding, there is a quantum mechanical electron exchange energy, which can be more simply explained that the sharing of electrons forms a bond.
Deciding whether there is bonding between two atoms is difficult and depends on the definition of what is a bond.\cite{nbo,dori}
For us, we intuitively say a bond has formed if the covalent radii of two atoms overlap.
These covalent radii are dictated by the element of the atom.

\begin{figure}[h]
    \centering
    \includegraphics[width=60ex]{figures/bond_formation.pdf}
    \caption{Defining a bond between atomic carbon and atomic oxygen}
    \label{fig:bond_definition}
\end{figure}

A simple example is given for carbon and oxygen (Fig.\ref{fig:bond_definition})
The covalent radii are imported from ASE, which are taken from ref.\cite{covradii}.
Regardless of whether the structure is stable, or if a quantum mechanical bond actually is formed, our strict definition is this.
With this definition of a bond, we can represent how molecules are connected to each other.

\subsection{Representing atoms as graphs}
Perhaps, as an imprinting of chemistry education, the graphical representation of molecules makes sense to chemists;
carbon can be bonded to 4 other atoms, whilst hydrogen can only be bonded to one.
Nodes are atomic nuclei, while edges are bonds.
Graphs only care about coordination however, and have no concern for position.
An example of this is shown in Fig.\ref{fig:molecular_graph}.

\begin{figure}[h]
    \centering
    \includegraphics[width=60ex]{figures/molecular_graph.pdf}
    \caption{Example molecular structure and graph representation for benzene}
    \label{fig:molecular_graph}
\end{figure}

There is an information loss cost involved with this graphical representation.
Stereoisomerism and chirality are not considered (Fig.\ref{fig:stereoisomers}).
However, what is lost is made up for by \textbf{rapid calculation}.
Comparing the similarity of structures is difficult and slow.
However, the comparison (or identification of a graph as an isomorph) can rapidly be deployed.
This is important for us as it lets us \textbf{find if a surface has a certain graphical representation that can undergo a reaction}.

\begin{figure}[h]
    \centering
    \includegraphics[width=60ex]{figures/stereoisomers.pdf}
    \caption{Example of a failure of graph representation, where isomorphic graphs are produced for cis/trans 1,2 dichloroethene}
    \label{fig:stereoisomers}
\end{figure}

\subsection{Graphical representation of surface chemistry}
\noindent



\section{Examples}

\section{Plan}

